Every synthesiser you will ever open is a variation on the same three-stage chain, with two
kinds of controller wired across it. Learn the chain once and every plugin becomes legible.

\subsection{The chain}\label{th05_sound_design-the-chain}

\begin{verbatim}
OSCILLATOR  ->  FILTER  ->  AMPLIFIER  ->  out
                  ^            ^
                  |            |
              ENVELOPE     ENVELOPE
                 LFO          LFO
\end{verbatim}

The \textbf{oscillator} makes a raw tone with a lot of harmonic content.

The \textbf{filter} removes some of those harmonics. This is where most of the character comes
from.

The \textbf{amplifier} controls the volume over time, which is what makes a sound a stab, a pluck
or a pad.

\textbf{Envelopes} and \textbf{LFOs} are not audible themselves. They are controllers that move the
other things around over time.

This is called \textbf{subtractive synthesis}, because you start with something harmonically rich
and carve away. It is the dominant method in dance music and it is what the TB-303, the
Juno-60 and roughly everything else in this folder are doing.

\subsection{Waveforms}\label{th05_sound_design-waveforms}

The oscillator\textquotesingle s shape determines which harmonics are present. The harmonics are integer
multiples of the fundamental frequency, and the amplitude pattern is what your ear reads as
timbre.

{\def\LTcaptype{none} % do not increment counter
\begin{longtable}[]{@{}>{\raggedright\arraybackslash}p{(\linewidth - 8\tabcolsep) * \real{0.0691}}>{\raggedright\arraybackslash}p{(\linewidth - 8\tabcolsep) * \real{0.1199}}>{\raggedright\arraybackslash}p{(\linewidth - 8\tabcolsep) * \real{0.1551}}>{\raggedright\arraybackslash}p{(\linewidth - 8\tabcolsep) * \real{0.6559}}@{}}
\toprule\noalign{}
Wave & Harmonics present & Amplitude falls off as & Sounds like \\
\midrule\noalign{}
\endhead
\bottomrule\noalign{}
\endlastfoot
Sine & fundamental only & n/a & pure, hollow, a whistle. The building block of sub bass \\
Triangle & odd only & 1/n² & soft, flute-like, slightly more body than a sine \\
Square & odd only & 1/n & hollow, woody, reedy. Classic for basses and leads \\
Sawtooth & all & 1/n & bright, buzzy, aggressive. The most harmonically rich, and the default for pads and supersaws \\
Noise & all, randomly & flat & hiss. Used for percussion, risers and texture \\
\end{longtable}
}

Two practical consequences.

A sine has nothing to filter, so a lowpass filter does nothing interesting to it. That is
why sub basses are usually sines with no filter movement at all.

A saw has everything to filter, so it responds enormously to filter movement. That is why
every filter sweep you have ever heard in dance music is a saw wave underneath.

A \textbf{pulse wave} is a square with an adjustable duty cycle. Modulating that duty cycle
(pulse width modulation, PWM) gives you the thick, shifting, slightly detuned sound of
1980s string machines. The Juno-60 that Larry Heard used for "Can You Feel It" is doing
exactly this.

\subsection{Filters}\label{th05_sound_design-filters}

A filter attenuates frequencies above or below a cutoff point.

\textbf{Lowpass (LPF)} passes low frequencies and removes highs. This is the one you will use 90\%
of the time. Closing a lowpass makes a sound darker, more distant, more muffled.

\textbf{Highpass (HPF)} does the opposite. Essential in mixing: highpass everything that is not
the kick and bass to clear out the low end.

\textbf{Bandpass} passes a band and removes everything either side. Telephone or radio effects,
and useful for making a sound sit in a narrow slot.

\textbf{Notch} removes a narrow band and passes the rest.

Two controls matter:

\textbf{Cutoff} is the corner frequency. Sweeping cutoff on a saw wave is the single most
recognisable gesture in electronic music.

\textbf{Resonance} boosts the frequencies right at the cutoff point. A little adds bite and
focus. A lot makes the filter ring and whistle, and at extreme settings the filter
self-oscillates and becomes a sine wave generator in its own right. Acid house is entirely
about high resonance.

\textbf{Slope} is how steeply the filter cuts, in decibels per octave. 12 dB/octave (2-pole) is
gentle. 24 dB/octave (4-pole) is the Moog ladder sound, tighter and more dramatic. The
TB-303 uses an unusual 18 dB/octave design, which is part of why acid basslines sound the way
they do and not like every other synth.

\subsection{Envelopes}\label{th05_sound_design-envelopes}

An envelope describes how a parameter changes over the life of a single note. The standard
form has four controls, universally abbreviated \textbf{ADSR}.

\textbf{Attack} is the time from the note starting to the peak. Zero attack means the sound
begins instantly, which is what you want for anything percussive. Long attack means it fades
in, which is what makes a pad a pad.

\textbf{Decay} is the time from that peak down to the sustain level.

\textbf{Sustain} is a \textbf{level, not a time}. It is how loud the sound stays while you keep holding
the key.

\begin{pitfallbox}
\textbf{Common mistake.} Reading sustain as a duration. Three of the four ADSR controls are
times and the fourth is not, which is why beginners set sustain to zero expecting silence
and get a pluck instead.
\end{pitfallbox}

\textbf{Release} is the time it takes to fade to silence after you let go.

An organ has full sustain and no decay: it stays until you release. A piano has zero sustain
and a long decay: it starts loud and fades even while you hold the key. A pluck is fast
attack, fast decay, zero sustain.

The single most useful trick in synthesis: route an envelope to the \textbf{filter cutoff} as
well as the volume. A short filter envelope on a saw wave gives you the "wow" at the front
of the note that makes a bass sound alive. Almost every good bass patch is doing this.

\subsection{LFOs}\label{th05_sound_design-lfos}

A low frequency oscillator is an oscillator running below hearing range, used as a
controller rather than as a sound. Typical rates are a fraction of a hertz up to about 20 Hz.

Route an LFO to:

\begin{itemize}
\tightlist
\item
  \textbf{pitch}, slowly and shallowly, for vibrato
\item
  \textbf{volume} for tremolo
\item
  \textbf{filter cutoff} for the wobbling, breathing quality on pads and, taken to extremes, for
  dubstep basses
\item
  \textbf{pan} for movement in the stereo field
\end{itemize}

In dance music LFOs are almost always \textbf{tempo-synced}, meaning their rate is locked to a
musical division (one cycle per bar, per beat, per 1/8) instead of a free-running frequency.
FL Studio does this with a tempo-sync button on every LFO.

\subsection{Other kinds of synthesis, briefly}\label{th05_sound_design-other-kinds-of-synthesis-briefly}

\textbf{FM synthesis} uses one oscillator to modulate the frequency of another at audio rate.
This generates complex, often inharmonic sidebands. It is brilliant for bells, metallic
sounds, electric pianos and gnarly basses, and it is much harder to control intuitively than
subtractive. In FL Studio: \textbf{Sytrus}.

\textbf{Additive synthesis} builds sounds by summing individual sine partials, giving you direct
control over each harmonic. Very powerful, very fiddly. In FL Studio: \textbf{Harmor} and
\textbf{Ogun}.

\textbf{Wavetable synthesis} scans through a table of stored single-cycle waveforms, so the
timbre morphs as you move through the table. This is the dominant modern method and the
reason Serum became ubiquitous. In FL Studio: the wavetable mode in \textbf{FLEX}, or third
party plugins like Serum and Vital.

\textbf{Physical modelling} simulates the physics of a real instrument rather than the sound. In
FL Studio: \textbf{Sakura} for strings, \textbf{Drumaxx} for drums.

\textbf{Granular synthesis} chops audio into tiny grains and reassembles them. Textural and
strange. In FL Studio: \textbf{Fruity Granulizer}.

Do not learn all of these. Learn subtractive properly and you can make 95\% of house.

\subsection{Sampling}\label{th05_sound_design-sampling}

Sampling is using recorded audio instead of a synthesised waveform, and it is at least half
of what dance music is made of.

\textbf{One-shots} are single hits: a kick, a clap, a stab. Load, trigger, done. Your drums will
mostly be one-shots.

\textbf{Loops} are pre-recorded bars of audio. Useful but dangerous, because a loop someone else
made carries their groove and their sound, and a track built entirely from loops sounds like
a track built entirely from loops.

\textbf{Chopping and slicing} means cutting a longer sample into pieces and replaying them in a
new order. This is the foundation of hip-hop, jungle, and all of Chicago footwork. In FL
Studio: \textbf{Slicex} and \textbf{Fruity Slicer}.

\textbf{Time stretching} changes the length of a sample without changing its pitch, so you can
fit a sample recorded at 100 BPM into a 124 BPM track. FL Studio does this automatically
when you enable stretching on an audio clip.

Two things to know about \textbf{vocal samples}, which are central to house. First, the classic
Chicago approach is a full soulful vocal over the whole track, gospel-derived, sung properly.
Second, the modern approach is the \textbf{vocal chop}: a short fragment of a vocal, pitched and
sequenced like an instrument. Both are house. They just come from different decades.

On the legal question, because it will come up: sampling a commercial record without
clearance is copyright infringement, and the fact that everyone did it in 1987 does not mean
you can now. For learning, sample whatever you like. Before you release anything, use
royalty-free sample packs, vocals you have licensed or recorded yourself, or public domain
material.

\subsection{The sounds you actually need for house}\label{th05_sound_design-the-sounds-you-actually-need-for-house}

\textbf{The kick.} A 909 kick is a sine-ish body with a fast downward pitch envelope and a click
at the front. You can build one from scratch in FL\textquotesingle s \textbf{Fruit Kick} or \textbf{BassDrum} plugin,
but you should start with the sampled \texttt{909\ Kick.wav} that ships with FL and only build your
own once you know what you are trying to make.

\textbf{The clap.} Several noise bursts a few milliseconds apart, band-passed around 1 kHz, with
a short reverb tail. That tiny stagger is why a clap sounds like several people and not one.
Use \texttt{909\ Clap.wav}.

\textbf{The hats.} Filtered noise plus metallic partials, gated very short for closed and longer
for open. Highpass everything below about 300 Hz. \texttt{909\ CH\ 1.wav} and \texttt{909\ OH.wav}.

\textbf{The acid bass.} A sawtooth or square through a resonant lowpass, with a filter envelope,
plus two things that are specific to the 303: \textbf{slide} (portamento between overlapping
notes) and \textbf{accent} (certain notes hit harder and open the filter more). Without slide and
accent you have a filtered bass, not an acid line. FL Studio ships \textbf{Transistor Bass},
which is a 303 emulation, and it is genuinely good. Load it, set resonance high, and turn
the cutoff knob while a 16th-note pattern loops. That is acid house.

\textbf{The chord stab.} A short, filtered, heavily enveloped chord played on the offbeat. Often
a sampled organ, a Rhodes, or a brass patch. Fast attack, short decay, no sustain. The
Chicago version is usually drier and more mechanical than the New York version.

\textbf{The pad.} Saw or PWM, slow attack, long release, lowpass filter mostly closed, plenty of
reverb, gentle LFO on the cutoff. This is the layer that gives deep house its atmosphere. In
FL Studio, \textbf{FLEX} has a lot of usable ones out of the box.

\textbf{The Rhodes / electric piano.} FL\textquotesingle s \textbf{FL Keys} does a passable one. Play seventh and ninth
chords on it and you are immediately in deep house territory.

\subsection{Effects are sound design too}\label{th05_sound_design-effects-are-sound-design-too}

\textbf{Reverb} puts a sound in a space. Short reverbs make things feel close and real, long ones
make them feel distant and dreamy. Use it on a \textbf{send}, not on every channel individually
(see \texttt{07\_mixing.md}).

\textbf{Delay} is discrete echoes. Tempo-synced to a dotted eighth is the classic dance music
setting. A ping-pong delay alternates between left and right and instantly widens a mono
sound.

\textbf{Chorus} duplicates the signal, detunes and delays the copy slightly, and mixes it back.
Makes anything thicker and wider. The Juno\textquotesingle s built-in chorus is a large part of why those
records sound the way they do.

\textbf{Distortion and saturation} add harmonics. A little on a kick or a bass adds weight and
makes it audible on small speakers. A lot is a creative choice.

\textbf{Bitcrushing and sample rate reduction} deliberately degrade the signal. Good for that
gritty, early-hardware texture.

\textbf{Filter automation on an effect send} is how you do the classic house breakdown, where
everything gets muffled and distant and then snaps back.

\subsection{How to actually learn this}\label{th05_sound_design-how-to-actually-learn-this}

Do not read synth manuals. Do this instead, once, for an hour:

Open \textbf{3x Osc} in FL Studio, which is deliberately primitive. Set one oscillator to a
sawtooth and turn the other two off. Play and hold a note. Now, one at a time and listening
the whole way:

\begin{enumerate}
\def\labelenumi{\arabic{enumi}.}
\tightlist
\item
  Sweep the filter cutoff from fully open to fully closed.
\item
  Put resonance to about 70\% and sweep again.
\item
  Set the volume envelope attack to maximum, then back to zero.
\item
  Set the decay short and sustain to zero. You now have a pluck.
\item
  Route the envelope to the filter and repeat step 4.
\item
  Add an LFO to the cutoff, sync it to 1/4, and listen.
\end{enumerate}

After that hour, every synth you open for the rest of your life will make sense. The
controls are always the same controls.
